\chapter{Theoretical Background}
\label{02_Theoretical_Background}
\thispagestyle{empty}
Before presenting the proposed system, it is necessary to introduce the theoretical background related to financial time series and their analysis using artificial intelligence methods. This thesis focuses on the application of AI techniques to financial data, which requires understanding both the properties of time series and the models used to process them.
\vspace{3mm}
The following subsections describe the main characteristics of financial time series, the role of feature engineering, and selected artificial intelligence methods used in algorithmic trading. The chapter also introduces ensemble methods and decision fusion, as well as the concept of agents and multi-agent systems, which form the basis of the proposed approach.

\subsection*{Financial Markets and Time Series Characteristics}


A time series is a collection of observations indexed by the date of each observation ~\cite{hamilton2020time}. Most commonly, it is a sequence of values recorded at successive, equally spaced points in time, forming discrete-time data ~\cite{bangert2021machine}. Examples of such time series include environmental readings (e.g., temperature, rainfall), business data (e.g., sales, web traffic), and financial metrics such as asset prices. Time series are analyzed to uncover patterns, trends, and relationships that evolve over time, enabling informed decisions and predictions ~\cite{zhang2024deep}.
\vspace{3mm} 
In financial markets, observations of an asset’s price taken at equally spaced time intervals constitute a financial time series ~\cite{zhang2024deep}. These observations are typically represented in the form of OHLCV data, which includes the open, high, low, and close prices, as well as the traded volume within a given time interval. Depending on the chosen timeframe, such data can be recorded at different resolutions, ranging from high-frequency (e.g., tick or minute-level) to lower-frequency intervals such as daily or weekly data.
\vspace{3mm} 
Beyond aggregated price data, financial markets also generate more detailed information, such as limit order book (LOB) data, which captures the supply and demand dynamics of the market. While OHLCV data provides a compact representation of price movements, LOB data offers a more granular view of market activity and can be used to derive additional features for analysis ~\cite{tsantekidis2020using}.
\vspace{3mm}
Financial time series exhibit several properties that make modeling and prediction challenging. Returns typically show weak autocorrelation, meaning that past price changes provide limited direct information about future movements, while volatility measures, such as absolute or squared returns, often exhibit significant persistence. In addition, financial data is non-stationary, with statistical properties such as mean and variance changing over time, and displays volatility clustering, where periods of high variability tend to persist. Return distributions are also heavy-tailed, meaning that extreme events occur more frequently than expected under normal assumptions. Finally, nonlinear dependencies and long memory effects further complicate modeling, motivating the use of more advanced data-driven approaches  ~\cite{sewell2011characterization}. These characteristics highlight the complexity of financial time series and the challenges associated with their modeling, which require careful feature design and appropriate modeling approaches.
% \vspace{3mm} while this is 
% - Financial Markets and Time Series Characteristics
% Market properties:
% non-stationarity
% noise
% volatility clustering
% Efficient Market Hypothesis (briefly)
% Stylized facts:
% fat tails
% autocorrelation (weak)
% Time series nature of financial data
% Challenges:
% regime changes
% low signal-to-noise ratio
\subsection*{Feature Engineering in Financial Data}
OHLCV data
Technical indicators:
trend, momentum, volatility
Sentiment data:
news, social media
Multi-timeframe features
Feature limitations

\subsection*{Artificial Intelligence Methods in Algorithmic Trading}

The growing availability of financial data and advances in computational
power have driven the widespread adoption of artificial intelligence
methods in algorithmic trading. These methods range from classical
statistical approaches to modern deep learning architectures, each
capturing different aspects of market dynamics. Sezer et
al.~\cite{sezer2020financial} surveyed 140 deep learning studies
published between 2005 and 2019 and found that LSTM and its variants
dominated as the architecture of choice, with stock price forecasting,
trend forecasting, and index forecasting together accounting for more
than 70\% of all publications. Zhang et al.~\cite{zhang2024deep}
extended this picture to 2022, confirming the continued prominence of
recurrent architectures while noting growing interest in Transformer-based
models. Sonkavde et al.~\cite{sonkavde2023forecasting} provide a broader
review encompassing both classical and deep learning methods alongside
traditional statistical techniques.

\paragraph{Classical and statistical approaches.}
Before machine learning became prevalent, rule-based systems encoding
explicit expert knowledge — such as moving-average crossovers or
momentum thresholds — were the primary tools for automated trading.
While interpretable and computationally lightweight, they do not adapt
to changing market conditions~\cite{sonkavde2023forecasting}. Statistical
models such as ARIMA and its extensions capture linear temporal
dependencies and serve as theoretically grounded baselines, but their
assumption of linearity and stationarity limits their effectiveness on
financial data where nonlinear dynamics and regime changes are the
norm~\cite{zhang2024deep, sonkavde2023forecasting}.

\paragraph{Machine learning methods.}
Machine learning methods learn patterns directly from data without
requiring a predefined functional form. Tree-based ensembles — most
notably Random Forest and XGBoost — have demonstrated competitive
performance across multiple benchmarks. Sonkavde et
al.~\cite{sonkavde2023forecasting} showed that combining such models with
LSTM outperformed any individual model, illustrating the benefit of
heterogeneous ensembles. Support vector machines and $k$-nearest
neighbours are also widely used, though they scale poorly to
high-dimensional feature spaces without careful preprocessing.

\paragraph{Deep learning for time series.}
Deep learning has become the dominant paradigm for financial forecasting
due to its ability to learn hierarchical representations without heavy
reliance on manual feature design~\cite{zhang2024deep}. LSTM networks
address the vanishing gradient problem of plain RNNs through gating
mechanisms that selectively retain or discard information over long
sequences, making them the most widely applied architecture in the
reviewed literature~\cite{sezer2020financial, zhang2024deep}. The GRU
offers a lighter alternative with comparable accuracy. Temporal
Convolutional Networks (TCNs) use dilated causal convolutions to capture
multi-scale local patterns efficiently, while hybrid CNN-LSTM
architectures combine local feature extraction with global temporal
modelling. Transformer models leverage self-attention to capture
long-range dependencies without sequential computation; efficient
variants such as Informer, Autoformer, and PatchTST have specifically
addressed the scalability challenges of applying Transformers to long
financial sequences~\cite{zhang2024deep}.

\paragraph{Reinforcement learning.}
Reinforcement learning frames trading as a sequential decision problem
in which an agent selects actions — buy, sell, or hold — and receives
rewards reflecting realised returns. Unlike supervised methods, RL agents
directly optimise cumulative trading performance, making the framework
naturally aligned with the profit-maximisation objective. Value-based
approaches such as Deep Q-Networks and policy gradient methods such as
Proximal Policy Optimisation are the most common formulations. Sezer et
al.~\cite{sezer2020financial} identify deep RL as one of the most
promising directions for future research, particularly for applications
requiring continuous adaptation to changing market conditions.

\paragraph{Evolutionary computation.}
Evolutionary methods such as genetic algorithms and neuroevolution
provide gradient-free optimisation applicable to strategy design, feature
selection, and hyperparameter tuning. In the financial deep learning
literature they have been used primarily as auxiliary optimisers rather
than standalone strategy generators~\cite{sonkavde2023forecasting,
sezer2020financial}, though neuroevolution approaches additionally offer
the ability to evolve network architectures alongside weights when
gradient information is unavailable or unreliable.

3.1 Classical approaches
rule-based strategies
statistical models (ARIMA, etc.)
3.2 Machine learning methods
regression, tree-based models
3.3 Deep learning for time series
LSTM, GRU, Transformers
temporal dependencies
3.4 Reinforcement learning
trading as decision process
reward-based learning
3.5 Evolutionary computation
genetic algorithms
neuroevolution
can
\subsection*{Ensemble Methods and Decision Fusion}
ensemble learning (bagging, boosting)
weighted voting
stacking / meta-learning
mixture of experts
\subsection*{Multi-Agent Systems}
5.1 Definition of agents
autonomy
decision-making entities
5.2 Multi-agent system types
centralized
decentralized
hierarchical
5.3 Cooperation vs competition
5.4 Multi-agent systems in finance (if literature exists)
2.1 Financial Markets & Data
OHLCV, indicators (SMA, RSI, MACD)
market efficiency, noise, non-stationarity
paper: The Relation between Price Changes and Trading Volume: A Survey
2.2 Artificial Intelligence in Trading
time-series forecasting
reinforcement learning in trading
sentiment analysis
2.3 Agents and Multi-Agent Systems
This chapter provides an overview of 
There is a great amount of articles regarding predicting time series.
There are also lots regarding the special type of time series — that is financial data. However, this space is never filled — with how extensive the approaches may be regarding to trading, it seems that the space will never be fully discovered. Some may say that this kind of assumption that we can leverage mathematics, and AI being a derivative from it, to successfully move in the market is impossible, since markets rely in a heavy way on market sentiment, and it might be really hard to include some crucial news into any agent's abilities — [TBD - is NLP analysis included]. Like broad geopolitical events, or news
Most of them include trading using only one method. 
While the concept of intelligent agents dates back to the 1990s (Russell & Norvig, 1995; Wooldridge & Jennings, 1995), multi-agent systems have recently gained renewed attention due to the emergence of Large Language Models (LLMs), which enable more flexible reasoning, communication, and tool use.
This paper follows the approach in which each agent implements a single method. An agent is understood as a computer system that exhibits the following properties~\cite{wooldridge1995intelligent}:
\begin{itemize}
    \item \textbf{autonomy}: operates without direct human intervention and has control over its actions and internal state;
    \item \textbf{social ability}: interacts with other agents (and possibly humans) via a communication mechanism;
    \item \textbf{reactivity}: perceives its environment and responds to changes in a timely manner;
    \item \textbf{pro-activeness}: exhibits goal-directed behaviour by taking initiative, not only reacting to the environment.
\end{itemize}
While the concept of intelligent agents dates back to the 1990s, multi-agent systems recently gained a leap in attention due to fast-paced Large Language Models (LLMs) adoption.
Multi-agents is not the most recent concept —— the definition of the agent [... year, definition]. The definition of multi-agent [.... year, definition]. Definitely the whole thing got very popular within last years - with the boom of Large Language Models. Most of the agent's  
